\section{Exhaustive squeeze and defected budgets}
\label{sec:exhaustive_squeeze_and_budgets}

In applications the domination records are rarely available on \(Y\)
itself.  More often one controls a finite piece of the problem---a
window, a truncation, a finite-dimensional subspace---and must account
separately for what the piece misses.  This section records the two
refinements of the confinement theorem needed in that situation.
Throughout, \(\IOL\) is an involutive object ledger with
\(\int_X\|\psim\|^2\,d\mu<\infty\).

\begin{definition}[Exhaustive moving ledger]
\label{def:duality_confinement:exhaustive-moving-ledger}
An \emph{exhaustive moving ledger} for \(\AX\) consists of, for each
\(n\geq1\), a real Hilbert space \(Y_n\), a bounded operator
\(\iota_n:Y_n\to Y\), a positive trace-class operator \(A_{X,n}\) on
\(Y_n\) and a positive trace-class operator \(T_n\) on \(Y\), such that
\[
  \AX\preceq\iota_nA_{X,n}\iota_n^*+T_n\qquad(n\geq1).
\]
It is \emph{vanishing-exhaustive} if there are positive trace-class
operators \(B_n\) on \(Y_n\) with \(A_{X,n}\preceq B_n\) and
\[
  \trace(\iota_nB_n\iota_n^*)+\trace T_n\to0\qquad(n\to\infty).
\]
\end{definition}

The operator \(A_{X,n}\) is the part of the ledger seen by the
\(n\)-th window, \(\iota_n\) transports it into \(Y\), and \(T_n\) is a
positive tail accounting for what the window misses.  Despite the
name, no exhaustion, monotonicity or convergence of the windows is
required; the name reflects the typical use, illustrated next.

\begin{theorem}[Exhaustive ledger squeeze]
\label{thm:duality_confinement:exhaustive-squeeze}
If \(\AX\) admits a vanishing-exhaustive moving ledger, then
\(\AX=0\) and \(\psim(x)=0\) for \(\mu\)-almost every \(x\).  If, in
addition, \(\psim\) is qualitatively separating, then \(Jx=x\) for
\(\mu\)-almost every \(x\).
\end{theorem}

\begin{proof}
Conjugation preserves the Loewner order: if \(A\preceq B\) on \(Y_n\),
then for \(h\in Y\),
\(\langle\iota_n(B-A)\iota_n^*h,h\rangle
=\langle(B-A)\iota_n^*h,\iota_n^*h\rangle\geq0\), so
\(\iota_nA\iota_n^*\preceq\iota_nB\iota_n^*\).  Hence
\[
  \AX\preceq\iota_nA_{X,n}\iota_n^*+T_n
  \preceq\iota_nB_n\iota_n^*+T_n=:B_n' .
\]
Each \(B_n'\) is positive and trace-class, because the trace class is
an ideal, and \(\trace B_n'\to0\) by hypothesis.
\Cref{thm:duality_confinement:master-theorem} applied to the sequence
\((B_n')\) gives the conclusions.
\end{proof}

\Cref{thm:duality_confinement:exhaustive-squeeze} is
\Cref{thm:duality_confinement:master-theorem} applied to the records
\(B_n'\); its content is bookkeeping, recording the data an
application typically computes.

\begin{example}[Finite-dimensional windows]
Let \(Y\) be separable, and let \(P_n\) be orthogonal projections onto
finite-dimensional subspaces \(Y_n\subseteq Y\) with \(P_n\to I\)
strongly, let \(\iota_n\) be the
inclusion of \(Y_n\), so that \(\iota_n\iota_n^*=P_n\), and put
\(Q_n=I-P_n\).  With \(\AX=V^*V\) as in
\Cref{rem:duality_confinement:analysis-operator}, the case \(t=1\) of
the operator inequality displayed below, applied to
\(V=VP_n+VQ_n\), gives
\[
  \AX\preceq2P_n\AX P_n+2Q_n\AX Q_n .
\]
So the compressions \(A_{X,n}=2\,\iota_n^*\AX\iota_n\) and the tails
\(T_n=2Q_n\AX Q_n\) form an exhaustive moving ledger.  The tails
satisfy \(\trace T_n=2\|VQ_n\|_{\mathrm{HS}}^2\to0\) by dominated
convergence, because \(V\) is Hilbert--Schmidt.  Whether the window
bounds \(B_n\) can be chosen with \(\trace(\iota_nB_n\iota_n^*)\to0\) is
the substantive question in any application.
\end{example}

The second refinement concerns a bound with two error terms of
different origin.  Such bounds come from the operator form of the
elementary inequality \((\alpha+\beta)^2\leq(1+t)\alpha^2+(1+t^{-1})\beta^2\):
for bounded operators \(V_1,V_2\) and \(t>0\),
\[
  (1+t)V_1^*V_1+(1+t^{-1})V_2^*V_2-(V_1+V_2)^*(V_1+V_2)
  =\bigl(\sqrt tV_1-V_2/\sqrt t\bigr)^*\bigl(\sqrt tV_1-V_2/\sqrt t\bigr)
  \succeq0 .
\]
So if the analysis operator of \Cref{rem:duality_confinement:analysis-operator}
splits as \(V=V_1+V_2\), then \(\AX\) is dominated, for every \(t>0\),
by \((1+t)V_1^*V_1+(1+t^{-1})V_2^*V_2\).

\begin{definition}[Defected obstruction budget]
\label{def:duality_confinement:defected-budget}
A \emph{defected obstruction budget} for \(\AX\) consists of positive
trace-class operators \(\Kminus\), \(E_{\mathrm{src}}\) and
\(E_{\mathrm{br}}\) on \(Y\) such that
\[
  \AX
  \preceq
  (1+t)\bigl(\Kminus+E_{\mathrm{src}}\bigr)+(1+t^{-1})E_{\mathrm{br}}
  \qquad\text{for every }t>0 .
\]
Here \(E_{\mathrm{src}}\) is a source defect and \(E_{\mathrm{br}}\) a
bridge defect.  A \emph{sequence} of defected obstruction budgets is a
sequence \((\Kminus_n,E_{\mathrm{src},n},E_{\mathrm{br},n})\), each a
defected obstruction budget for the same operator \(\AX\).
\end{definition}

The operators \(\Kminus\) and \(E_{\mathrm{src}}\) enter only through
their sum; they are kept apart because in applications they have
different origins, a computed main term and a defect in its source.

The choice \(t=1\) gives the unweighted bound
\(2(\Kminus+E_{\mathrm{src}})+2E_{\mathrm{br}}\); the next proposition
optimizes over \(t\).

\begin{proposition}[Optimized scalar trace budget]
\label{prop:duality_confinement:optimized-trace-budget}
Let \((\Kminus,E_{\mathrm{src}},E_{\mathrm{br}})\) be a defected
obstruction budget for \(\AX\), and put
\(a=\trace(\Kminus+E_{\mathrm{src}})\) and \(b=\trace E_{\mathrm{br}}\).
Then \(a,b\geq0\),
\[
  \inf_{t>0}\bigl[(1+t)a+(1+t^{-1})b\bigr]
  =\bigl(\sqrt a+\sqrt b\bigr)^2,
\]
the infimum being attained at \(t=\sqrt{b/a}\) when \(a,b>0\), and
\[
  \trace\AX\leq\bigl(\sqrt a+\sqrt b\bigr)^2 .
\]
Consequently, if \((\Kminus_n,E_{\mathrm{src},n},E_{\mathrm{br},n})\)
is a sequence of defected obstruction budgets for \(\AX\) with
\(a_n=\trace(\Kminus_n+E_{\mathrm{src},n})\to0\) and
\(b_n=\trace E_{\mathrm{br},n}\to0\), then all conclusions of
\Cref{thm:duality_confinement:master-theorem} hold.
\end{proposition}

\begin{proof}
The traces \(a\) and \(b\) are nonnegative by (T3).  For \(t>0\),
\[
  (1+t)a+(1+t^{-1})b-\bigl(\sqrt a+\sqrt b\bigr)^2
  =ta+\frac bt-2\sqrt{ab}
  =\Bigl(\sqrt{ta}-\sqrt{b/t}\Bigr)^2\geq0,
\]
with equality when \(ta=b/t\), that is, at \(t=\sqrt{b/a}\) if
\(a,b>0\).  If \(a=0\) the budget equals \(b+b/t\), which decreases to
\(b=(\sqrt a+\sqrt b)^2\) as \(t\to\infty\); if \(b=0\) it equals
\(a+ta\), which decreases to \(a\) as \(t\to0\).  This proves the
formula for the infimum.  For each \(t>0\), (T4) and linearity of the
trace give \(\trace\AX\leq(1+t)a+(1+t^{-1})b\); taking the infimum
gives \(\trace\AX\leq(\sqrt a+\sqrt b)^2\).  For a sequence of budgets,
\(0\leq\trace\AX\leq(\sqrt{a_n}+\sqrt{b_n})^2\to0\), so
\(\trace\AX=0\), hence \(\AX=0\) by (T3), and the remaining
conclusions follow from
\Cref{thm:duality_confinement:separation-confinement}.
\end{proof}

When \(\AX=V^*V\) and \(V=V_1+V_2\) with
\(\Kminus+E_{\mathrm{src}}=V_1^*V_1\) and \(E_{\mathrm{br}}=V_2^*V_2\),
as in the operator inequality above, the bound
\(\trace\AX\leq(\sqrt a+\sqrt b)^2\) is the triangle inequality
\(\|V\|_{\mathrm{HS}}\leq\|V_1\|_{\mathrm{HS}}+\|V_2\|_{\mathrm{HS}}\),
and the optimization over \(t\) is a standard proof of it.  The
proposition applies more generally, to any family of bounds of the
displayed form.

For a single budget, combining the proposition with
\Cref{thm:duality_confinement:quantitative-confinement} gives the
explicit bound
\(\mu^*\{x:\dist(x,\Fix(J))\geq\varepsilon\}\leq
(\sqrt a+\sqrt b)^2/m(\varepsilon)^2\) whenever \(\psim\) is
quantitatively separating with modulus \(m\).
